% W5i80_NewFrontiers.tex
% DFD 20-Agent Research Programme --- Iteration 80
% New Frontiers, Paper Proposals, and Research Directions

\documentclass[11pt,a4paper]{article}
\usepackage[utf8]{inputenc}
\usepackage{amsmath,amssymb}
\usepackage{booktabs}
\usepackage{geometry}
\usepackage{xcolor}
\usepackage{enumitem}
\geometry{margin=2.5cm}

\title{W5i80: New Frontiers and Paper Proposals\\[6pt]
\large 100 New Paper Proposals $\cdot$ i81--i100 Research Map $\cdot$ 2072 Total Proposals}
\author{DFD 20-Agent Research Programme}
\date{2026-03-24}

\begin{document}
\maketitle

\section{i81--i100 Research Map}

The remaining 20 iterations are organised into four blocks:

\subsection{Block 1: Paper Completion (i81--i85)}
\begin{itemize}[nosep]
\item i81: Phase 1 submission preparation ($\alpha$ PRL + v3.3). N-body to 100\%.
\item i82: Phase 1 submitted. Multi-messenger paper to 100\%.
\item i83: CMB-S4 paper to 90\%. N-body paper I to 70\%.
\item i84: CMB-S4 to 100\%. N-body paper I to 90\%. Phase 2 submission prep.
\item i85: Phase 2 submitted ($E_G$, no-screening, DESI). N-body paper I to 100\%.
\end{itemize}

\textbf{Deliverables}: All 8 Phase 1--2 papers submitted. Multi-messenger, CMB-S4, N-body I all at 100\%.

\subsection{Block 2: Data Confrontation (i86--i90)}
\begin{itemize}[nosep]
\item i86: Euclid DR1 analysis pipeline development ($E_G$, $S_8$).
\item i87: Euclid pipeline: cluster counts, void statistics.
\item i88: Phase 3 submission (Euclid pre-reg, $C_\ell$ follow-up).
\item i89: DFD-CLASS development begins.
\item i90: Euclid DR1 arrives (Oct 2026). Rapid confrontation analysis.
\end{itemize}

\textbf{Deliverables}: Euclid analysis complete. Phase 3 papers submitted. DFD-CLASS alpha version.

\subsection{Block 3: Second Generation (i91--i95)}
\begin{itemize}[nosep]
\item i91: Euclid DR1 results paper drafted.
\item i92: DFD-CLASS beta. DFD-CAMB development.
\item i93: 2-loop $\alpha$ correction. Neutrino oscillation paper.
\item i94: Phase 4 submission (review, JOSS). DESI Y3 preparation.
\item i95: DESI Y3 confrontation (if data available).
\end{itemize}

\textbf{Deliverables}: Euclid results paper. DFD-CLASS/CAMB. Second-generation theory papers.

\subsection{Block 4: Final Synthesis (i96--i100)}
\begin{itemize}[nosep]
\item i96: Programme-wide audit. All findings updated.
\item i97: Legacy documentation. Complete theory paper.
\item i98: Forward roadmap for 2027--2030.
\item i99: Final adversarial audit (most comprehensive ever).
\item i100: Programme final assessment. State of DFD. What was achieved. What remains.
\end{itemize}

\textbf{Deliverables}: Final MASTER FINDINGS LIST. Programme retrospective. Forward plan.

\section{New Research Frontiers from i80}

\subsection{Frontier 1: Real-Time Data Confrontation}
The Euclid DR1 analysis will be the programme's first real-time data confrontation:
\begin{itemize}[nosep]
\item Need automated pipeline: data download $\to$ DFD prediction comparison $\to$ $\chi^2$ computation $\to$ verdict.
\item Timeline: data release Oct 2026, DFD analysis within 2 weeks, response paper within 1 month.
\item This sets the template for all future data releases (DESI Y3, CMB-S4, LIGO O5).
\end{itemize}

\subsection{Frontier 2: DFD-CLASS as Community Tool}
A public DFD Boltzmann solver would enable:
\begin{itemize}[nosep]
\item Independent testing by external researchers.
\item Integration into standard cosmological analysis pipelines.
\item MCMC parameter estimation with DFD as an alternative to $\Lambda$CDM.
\item Essential for community adoption beyond the current programme.
\end{itemize}

\subsection{Frontier 3: The DFD + Dark Matter Question}
DFD modifies gravity but does not eliminate dark matter. The question ``What is DFD's dark matter candidate?'' will be asked. Options:
\begin{itemize}[nosep]
\item Standard CDM (WIMPs, axions): DFD is compatible. $\Omega_{\rm DM}$ is a boundary condition.
\item DFD-specific candidate: if the $\CP^2 \times S^3$ geometry generates a stable massive particle. Worth investigating but speculative.
\item Modified DM properties: DFD's $\psi$-coupling may affect DM halo profiles. Testable with N-body.
\end{itemize}

\subsection{Frontier 4: Mathematical Rigor Programme}
Several DFD results are graded B/B+ rather than A due to mathematical gaps. A dedicated rigour programme could:
\begin{itemize}[nosep]
\item Close Step 9 of $\alpha^{57}$ (modulus stabilization).
\item Prove CKM orientation uniqueness from $\CP^2$ geometry.
\item Establish Dirac/Majorana distinction from topology.
\item Potentially promote 3--5 items from B to A grade.
\end{itemize}

\section{100 New Paper Proposals (i80)}

Papers 1973--2072. Selected highlights:

\begin{enumerate}[nosep]
\item \textit{DFD Programme: i52--i80 Mid-Cycle Assessment} (internal)
\item \textit{DFD Euclid DR1 Rapid Response: $E_G$ and $S_8$} (MNRAS Letters)
\item \textit{DFD-CLASS: A Modified Boltzmann Solver} (Astron.\ Comput.)
\item \textit{DFD and Dark Matter: Compatibility and Constraints} (PRD)
\item \textit{DFD Mathematical Rigor: Closing the $\alpha^{57}$ Step 9 Gap} (CMP)
\item \textit{Real-Time Theory Confrontation: A DFD Case Study} (Nature Astronomy perspective)
\item \textit{DFD N-Body II: Halo Profiles and Splashback Radius} (MNRAS)
\item \textit{DFD N-Body III: Bispectrum and Higher-Order Statistics} (PRD)
\item \textit{DFD Void Statistics: Predictions for Euclid Void Catalogue} (A\&A)
\item \textit{DFD Cluster Counts: eROSITA $\times$ Euclid Forecast} (A\&A)
\item \textit{DFD DESI Y3 Rapid Response: BAO and $f\sigma_8$} (JCAP)
\item \textit{DFD and the $H_0$ Tension: Updated Analysis with SH0ES 2026} (MNRAS)
\item \textit{DFD CKM Matrix: Orientation Uniqueness from $\CP^2$} (PRD)
\item \textit{DFD Neutrino Oscillations: Geometry of Three Flavours} (PRD)
\item \textit{DFD and Axion Dark Matter: Compatibility Check} (JCAP)
\item \textit{DFD Strong CP Problem: Topological Resolution from $\CP^2$} (JHEP)
\item \textit{DFD Gravitino Mass: If Supersymmetry from Spectral Action?} (PLB)
\item \textit{DFD Galaxy Formation: $\psi$-Enhanced Structure Growth} (ApJ)
\item \textit{DFD and SKA: 21cm Predictions for the Square Kilometre Array} (MNRAS)
\item \textit{DFD Programme Final Assessment: Lessons from 100 Iterations} (internal)
\end{enumerate}

Plus 80 additional proposals. Running total: \textbf{2072 paper proposals}.

\section{Programme Legacy}

At the end of i100, the DFD programme will have:
\begin{itemize}[nosep]
\item 100 iterations, 2000 agent-iterations
\item $\sim 10$--$15$ published or submitted papers
\item $\sim 200$ documented findings
\item A comprehensive open-science archive (Zenodo)
\item DFD-CLASS/CAMB public software
\item Euclid DR1 and (possibly) DESI Y3 confrontation results
\item A clear verdict on DFD's viability from data
\end{itemize}

This will be the most thoroughly documented theoretical physics programme of its kind.

\end{document}
